\chapter[Chapter \thechapter]{}
\lettrine[ante=`]{B}{unter} !'

\zz
<My lord?>

\zz
Wimsey tapped with his fingers a letter he had just received.

\zz
<Do you feel at your brightest and most truly fascinating? Does a livelier iris, winter weather notwithstanding, shine upon the burnished Bunter? Have you got that sort of conquering feeling? The Don Juan touch, so to speak?>

Bunter, balancing the breakfast tray on his fingers, coughed deprecatingly.

<You have a good, upstanding, impressive figure, if I may say so,> pursued Wimsey, <a bold and roving eye when off duty, a ready tongue, Bunter—and, I am persuaded, you have a way with you. What more should any cook or house-parlourmaid want?>

<I am always happy,> replied Bunter, <to exert myself to the best of my capacity in your lordship's service.>

<I am aware of it,> admitted his lordship. <Again and again I say to myself, Wimsey, this cannot last. One of these days this worthy man will cast off the yoke of servitude and settle down in a pub. or something, but nothing happens. Still, morning by morning, my coffee is brought, my bath is prepared, my razor laid out, my ties and socks sorted and my bacon and eggs brought to me in a lordly dish. No matter. This time I demand a more perilous devotion—perilous for us both, my Bunter, for if you were to be carried away, a helpless martyr to matrimony, who then would bring my coffee, prepare my bath, lay out my razor and perform all those other sacrificial rites? And yet\longdash>

<Who is the party, my lord?>

<There are two of them, Bunter, two ladies lived in a bower, Binnorie, O Binnorie! The parlourmaid you have seen. Her name is Hannah Westlock. A woman in her thirties, I fancy, and not ill-favored. The other, the cook—I cannot lisp the tender syllables of her name, for I do not know it, but doubtless it is Gertrude, Cecily, Magdalen, Margaret, Rosalys or some other sweet symphonious sound—a fine woman, Bunter, on the mature side, perhaps, but none the worse for that.>

<Certainly not, my lord. If I may say so, the woman of ripe years and queenly figure is frequently more susceptible to delicate attentions than the giddy and thoughtless young beauty.>

<True. Let us suppose, Bunter, that you were to be the bearer of a courteous missive to one Mr~Norman Urquhart of Woburn Square. Could you, in the short space of time at your disposal, insinuate yourself, snakelike, as it were, into the bosom of the household?>

<If you desire it, my lord, I will endeavour to insinuate myself to your lordship's satisfaction.>

<Noble fellow. In case of an action for breach, or any consequence of that description, the charges will, of course, be borne by the management.>

<I am obliged to your lordship. When would your lordship wish me to commence?>

<As soon as I have written a note to Mr~Urquhart. I will ring.>

<Very good, my lord.>

Wimsey moved over to the writing-desk. After a few moments he looked up, a little peevishly.

<Bunter, I have a sensation of being hovered over. I do not like it. It is unusual and it unnerves me. I implore you not to hover. Is the proposition distasteful, or do you want me to get a new hat? What is troubling your conscience?>

<I beg your lordship's pardon. It had occurred to my mind to ask your lordship, with every respect\longdash>

<Oh, God, Bunter—don't break it gently. I can't bear it. Stab and end the creature—to the heft! What is it?>

<I wished to ask you, my lord, whether your lordship thought of making any changes in your establishment?>

Wimsey laid down his pen and stared at the man.

<Changes, Bunter? When I have just so eloquently expressed to you my undying attachment to the loved routine of coffee, bath, razor, socks, eggs and bacon and the old, familiar faces? You're not giving me warning, are you?>

<No, indeed, my lord. I should be very sorry to leave your lordship's service. But I had thought it possible that, if your lordship was about to contract new ties\longdash>

<I \textit{knew} it was something in the haberdashery line! By all means, Bunter, if you think it necessary. Had you any particular pattern in mind?>

<Your lordship misunderstands me. I referred to domestic ties, my lord. Sometimes, when a gentleman reorganizes his household on a matrimonial basis, the lady may prefer to have a voice in the selection of the gentleman's personal attendant, in which case\longdash>

<Bunter!> said Wimsey, considerably startled, <may I ask where you have contracted these ideas?>

<I ventured to draw an inference, my lord.>

<This comes of training people to be detectives. Have I been nourishing a sleuth-hound on my own hearth-stone? May I ask if you have gone so far as to give a name to the lady?>

<Yes, my lord.>

There was a pause.

<Well?> said Wimsey, in a rather subdued tone, <what about it, Bunter?>

<A very agreeable lady, if I may say so, my lord.>

<It strikes you that way, does it? The circumstances are unusual, of course.>

<Yes, my lord. I might perhaps make so bold as to call them romantic.>

<You may make so bold as to call them damnable, Bunter.>

<Yes, my lord,> said Bunter, in a tone of sympathy.

<You won't desert the ship, Bunter?>

<Not on any account, my lord.>

<Then don't come frightening me again. My nerves are not what they were. Here is the note. Take it round and do your best.>

<Very good, my lord.>

<Oh, and, Bunter.>

<My lord?>

<It seems that I am being obvious. I have no wish to be anything of the kind. If you see me being obvious, will you drop me a hint?>

<Certainly, my lord.>

Bunter faded gently out, and Wimsey stepped anxiously to the mirror.

<I can't see anything,> he said to himself. <No lily on my cheek with anguish moist and fever-dew. I suppose, though, it's hopeless to try and deceive Bunter. Never mind. Business must come first. I've stopped one, two, three, four earths. What next? How about this fellow Vaughan?>

\divider

When Wimsey had any researches to do in Bohemia, it was his custom to enlist the help of Miss~Marjorie Phelps. She made figurines in porcelain for a living, and was therefore usually to be found either in her studio or in some one else's studio. A telephone-call at 10 \textsc{\textsc{a.m.}} would probably catch her scrambling eggs over her own gas-stove. It was true that there had been passages, about the time of the Bellona Club affair\footnote{See \textit{The Unpleasantness at the Bellona Club,} published 1928.}, between her and Lord~Peter which made it a little embarrassing and unkind to bring her in on the subject of Harriet Vane, but with so little time in which to pick and choose his tools, Wimsey was past worrying about gentlemanly scruples. He put the call through and was relieved to hear an answering <Hullo!>

<Hullo, Marjorie! This is Peter Wimsey. How goes it?>

<Oh, fine, thanks. Glad to hear your melodious voice again. What can I do for the Lord High Investigator?>

<Do you know one Vaughan, who is mixed up in the Philip Boyes murder mystery?>

<Oh, Peter! Are you on to that? How gorgeous! Which side are you taking?>

<For the defence.>

<Hurray!>

<Why this pomp of jubilee?>

<Well, it's much more exciting and difficult, isn't it?>

<I'm afraid it is. Do you know Miss~Vane, by the way?>

<Yes and no. I've seen her with the Boyes-Vaughan crowd.>

<Like her?>

<So-so.>

<Like him? Boyes, I mean?>

<Never stirred a heart-beat.>

<I said, did you like him?>

<One didn't. One either fell for him or not. He wasn't the merry bright-eyed pal of the period, you know.>

<Oh! What's Vaughan?>

<Hanger-on.>

<Oh?>

<House-dog. Nothing must interfere with the expansion of my friend the genius. That sort.>

<Oh!>

<Don't keep saying <Oh!> Do you want to meet the man Vaughan?>

<If it's not too much trouble.>

<Well, turn up tonight with a taxi and we'll go the rounds. We're certain to drop across him somewhere. Also the rival gang, if you want them—Harriet Vane's supporters.>

<Those girls who gave evidence?>

<Yes. You'll like Eiluned Price, I think. She scorns everything in trousers, but she's a good friend at a pinch.>

<I'll come, Marjorie. Will you dine with me?>

<Peter, I'd adore to, but I don't think I will. I've got an awful lot to do.>

<Right-ho! I'll roll round about nine, then.>

Accordingly, at 9 o'clock, Wimsey found himself in a taxi with Marjorie Phelps, headed for a round of the studios.

<I've been doing some intensive telephoning,> said Marjorie, <and I think we shall find him at the Kropotkys'. They are pro-Boyes, Bolshevik and musical, and their drinks are bad, but their Russian tea is safe. Does the taxi wait?>

<Yes; it sounds as if we might want to beat a retreat.>

<Well, it's nice to be rich. It's down the court here, on the right, over the Petrovitchs' stable. Better let me grope first.>

They stumbled up a narrow and encumbered stair, at the top of which a fine confused noise of a piano, strings and the clashing of kitchen utensils announced that some sort of entertainment was in progress.

Marjorie hammered loudly on a door, and, without waiting for an answer, flung it open. Wimsey, entering on her heels, was struck in the face, as by an open hand, by a thick muffling wave of heat, sound, smoke and the smell of frying.

It was a very small room, dimly lit by a single electric bulb, smothered in a lantern of painted glass, and it was packed to suffocation with people, whose silk legs, bare arms and pallid faces loomed at him like glow-worms out of the obscurity. Coiling wreaths of tobacco-smoke swam slowly to and fro in the midst. In one corner an anthracite stove, glowing red and mephitical, vied with a roaring gas-oven in another corner to raise the atmosphere to roasting-pitch. On the stove stood a vast and steaming kettle; on a side-table stood a vast and steaming samovar; over the gas, a dim figure stood turning sausages in a pan with a fork, while an assistant attended to something in the oven, which Wimsey, whose nose was selective, identified among the other fragrant elements in this compound atmosphere, and identified rightly, as kippers. At the piano, which stood just inside the door, a young man with bushy red hair was playing something of a Czecho-Slovakian flavour, to a violin obligato by an extremely loose-jointed person of indeterminate sex in a Fair-Isle jumper. Nobody looked round at their entrance. Marjorie picked her way over the scattered limbs on the floor and, selecting a lean young woman in red, bawled into her ear. The young woman nodded and beckoned to Wimsey. He negotiated a passage and was introduced to the lean woman by the simple formula: <Here's Peter—this is Nina Kropotky.>

<So pleased,> shouted Madame Kropotky through the clamour. <Sit by me. Vanya will get you something to drink. It is beautiful, yes? That is Stanislas—such a genius—his new work on the Piccadilly Tube Station—great, n'est-ce pas? Five days he was continually travelling upon the escalator to absorb the tone-values.>

<Colossal!> yelled Wimsey.

<So—you think? Ah! you can appreciate! You understand it is really for the big orchestra. On the piano it is nothing. It needs the brass, the effects, the timpani—b'rrrrrrr! So! But one seizes the form, the outline! Ah! it finishes! Superb! Magnificent!>

The enormous clatter ceased. The pianist mopped his face and glared haggardly round. The violinist put down its instrument and stood up, revealing itself, by its legs, to be female. The room exploded into conversation. Madame Kropotky leapt over her seated guests and embraced the perspiring Stanislas on both cheeks. The frying-pan was lifted from the stove in a fusillade of spitting fat, a shriek went up for <Vanya!> and presently a cadaverous face was pushed down to Wimsey's, and a deep guttural voice barked at him: <What will you drink?> while simultaneously a plate of kippers came hovering perilously over his shoulder.

<Thanks,> said Wimsey, <I have just dined—just \textit{dined},< he roared despairingly, >full up, complet!>

Marjorie came to the rescue with a shriller voice and more determined refusal.

<Take those dreadful things away, Vanya. They make me sick. Give us some tea, tea, tea!>

<Tea!> echoed the cadaverous man, <they want tea! What do you think of Stanislas' tone-poem? Strong, modern, eh? The soul of rebellion in the crowd—the clash, the revolt at the heart of the machinery. It gives the bourgeois something to think of, oh, yes!>

<Bah!> said a voice in Wimsey's ear, as the cadaverous man turned away, <it is nothing. Bourgeois music. Programme music. Pretty!—You should hear Vrilovitch's <Ecstasy on the letter Z.> That is pure vibration with no antiquated pattern in it. Stanislas—he thinks much of himself, but it is old as the hills—you can sense the resolution at the back of all his discords. Mere harmony in camouflage. Nothing in it. But he takes them all in because he has red hair and reveals his bony structure.>

The speaker certainly did not err along these lines, for he was as bald and round as a billiard-ball. Wimsey replied soothingly:

<Well, what can you do with the wretched and antiquated instruments of our orchestra? A diatonic scale, bah! Thirteen miserable, bourgeois semi-tones, pooh! To express the infinite complexity of modern emotion, you need a scale of thirty-two notes to the octave.>

<But why cling to the octave?> said the fat man. <Till you can cast away the octave and its sentimental associations, you walk in fetters of convention.>

<That's the spirit!> said Wimsey. <I would dispense with all definite notes. After all, the cat does not need them for his midnight melodies, powerful and expressive as they are. The love-hunger of the stallion takes no account of octave or interval in giving forth the cry of passion. It is only man, trammelled by a stultifying convention—Oh, hullo, Marjorie, sorry—what is it?>

<Come and talk to Ryland Vaughan,> said Marjorie. <I have told him you are a tremendous admirer of Philip Boyes' books. Have you read them?>

<Some of them. But I think I'm getting light-headed.>

<You'll feel worse in an hour or so. So you'd better come now.> She steered him to a remote spot near the gas-oven, where an extremely elongated man was sitting curled up on a floor cushion, eating caviare out of a jar with a pickle-fork. He greeted Wimsey with a sort of lugubrious enthusiasm.

<Hell of a place,> he said, <hell of a business altogether. This stove's too hot. Have a drink. What the devil else can one do? I come here, because Philip used to come here. Habit, you know. I hate it, but there's nowhere else to go.>

<You knew him very well, of course,> said Wimsey, seating himself in a waste-paper basket, and wishing he was wearing a bathing-suit.

<I was his only real friend,> said Ryland Vaughan, mournfully. <All the rest only cared to pick his brains. Apes! parrots! all the bloody lot of them.>

<I've read his books and thought them very fine,> said Wimsey, with some sincerity. <But he seemed to me an unhappy soul.>

<Nobody understood him,> said Vaughan. <They called him difficult—who wouldn't be difficult with so much to fight against? They sucked the blood out of him, and his damned thieves of publishers took every blasted coin they could lay their hands on. And then that bitch of a woman poisoned him. My God, what a life!>

<Yes, but what made her do it—if she did do it?>

<Oh, she did it all right. Sheer, beastly spite and jealousy, that's all there was to it. Just because she couldn't write anything but tripe herself. Harriet Vane's got the bug all these damned women have got—fancy they can do things. They hate a man and they hate his work. You'd think it would have been enough for her to help and look after a genius like Phil, wouldn't you? Why, damn it, he used to ask her advice about his work, her advice, good lord!>

<Did he take it?>

<Take it? She wouldn't give it. Told him she never gave opinions on other authors' work. \textit{Other} authors! The impudence of it! Of course she was out of things among us all, but why couldn't she realise the difference between her mind and his? Of course it was hopeless from the start for Philip to get entangled with that kind of woman. Genius must be served, not argued with. I warned him at the time, but he was infatuated. And then, to want to marry her\longdash>

<Why did he?> asked Wimsey.

<Remains of parsonical upbringing, I suppose. It was really pitiful. Besides, I think that fellow Urquhart did a lot of mischief. Sleek family lawyer—d'you know him?>

<No.>

<He got hold of him—put up to it by the family, I imagine. I saw the influence creeping over Phil long before the real trouble began. Perhaps it's a good thing he's dead. It would have been ghastly to watch him turn conventional and settle down.>

<When did this cousin start getting hold of him, then?>

<Oh—about two years ago—a little more, perhaps. Asked him to dinner and that sort of thing. The minute I saw him I knew he was out to ruin Philip, body and soul. What he wanted—what Phil wanted, I mean—was freedom and room to turn about in, but what with the woman and the cousin and the father in the background—oh, well! It's no use crying about it now. His work is left, and that's the best part of him. He's left me that to look after, at least. Harriet Vane didn't get her finger in that pie, after all.>

<I'm sure it's absolutely safe in your hands,> said Wimsey.

<But when one thinks what there might have been,> said Vaughan, turning his blood-shot eyes miserably on Lord~Peter, <it's enough to make one cut one's throat, isn't it?>

Wimsey expressed agreement.

<By the way,> he said, <you were with him all that last day, till he went to his cousin's. You don't think he had anything on him in the way of—poison or anything? I don't want to seem unkind—but he was unhappy—it would be rotten to think that he\longdash>

<No,> said Vaughan, <no. That I'll swear he never did. He would have told me—he trusted me in those last days. I shared all his thoughts. He was miserably hurt by that damned woman, but he wouldn't have gone without telling me or saying good-bye. And besides—he wouldn't have chosen that way. Why should he? I could have given him\longdash>

He checked himself, and glanced at Wimsey, but, seeing nothing in his face beyond sympathetic attention, went on:

<I remember talking to him about drugs. Hyoscine—veronal—all that sort of thing. He said, <If ever I want to go out, Ryland, you'll show me the way.> And I would have—if he'd really wanted it. But arsenic! Philip, who loved beauty so much—do you think he would have chosen arsenic?—the suburban poisoner's outfit? That's absolutely impossible.>

<It's not an agreeable sort of thing to take, certainly,> said Wimsey.

<Look here,> said Vaughan, hoarsely and impressively—he had been putting a constant succession of brandies on top of the caviare, and was beginning to lose his reserve—<Look here! See this!> He pulled a small bottle from his breast-pocket. <That's waiting, till I've finished editing Phil's books. It's a comfort to have it there to look at, you know. Peaceful. Go out through the ivory gate—that's classical—they brought me up on the classics. These people would laugh at a fellow, but you needn't tell them I said it—funny, the way it sticks—'tendebantque manus ripae ulterioris amore, ulterioris amore'—what's that bit about the souls thronging thick as leaves in Vallombrosa—no, that's Milton—'amorioris ultore'—ultoriore—damn it—poor Phil!>

Here Mr~Vaughan burst into tears and patted the little bottle.

Wimsey, whose head and ears were thumping as though he were sitting in an engine-room, got up softly and withdrew. Somebody had begun a Hungarian song and the stove was white-hot. He made signals of distress to Marjorie, who was sitting in a corner with a group of men. One of them appeared to be reading his own poems with his mouth nearly in her ear, and another was sketching something on the back of an envelope, to the accompaniment of yelps of merriment from the rest. The noise they made disconcerted the singer, who stopped in the middle of a bar, and cried angrily:

<Ach! this noise! these interruptions! they are intolerable! I lose myself! Stop! I begin all over again, from the beginning.>

Marjorie sprang up, apologising.

<I'm a brute—I'm not keeping your menagerie in order, Nina—we're being perfect nuisances. Forgive me, Marya, I'm in a bad temper. I'd better pick up Peter and toddle away. Come and sing to me another day, darling, when I'm feeling better and there is more room for my feelings to expand. Good-night, Nina—we've enjoyed it frightfully—and, Boris, that poem's the best thing you've done, only I couldn't hear it properly. Peter, tell them what a rotten mood I'm in tonight and take me home.>

<That's right,> said Wimsey, <nervy, you know—bad effect on the manners and so on.>

<Manners,> said a bearded gentleman suddenly and loudly, <are for the bourgeois.>

<Quite right,> said Wimsey. <Beastly bad form, and gives you repressions in the what-not. Come on, Marjorie, or we shall all be getting polite.>

<I begin again,> said the singer, <from the beginning.>

<Whew!> said Wimsey, on the staircase.

<Yes, I know. I think I'm a perfect martyr to put up with it. Anyway, you've seen Vaughan. Nice dopey specimen, isn't he?>

<Yes, but I don't think he murdered Philip Boyes, do you? I had to see him to make sure. Where do we go next?>

<We'll try Joey Trimbles. That's the stronghold of the opposition show.>

Joey Trimbles occupied a studio over a mews. Here there was the same crowd, the same smoke, more kippers, still more drinks and still more heat and conversation. In addition there was a blaze of electric light, a gramophone, five dogs and a strong smell of oil-paints. Sylvia Marriott was expected. Wimsey found himself involved in a discussion of free love, D. H. Lawrence, the prurience of prudery and the immoral significance of long skirts. In time, however, he was rescued by the arrival of a masculine-looking middle-aged woman with a sinister smile and a pack of cards, who proceeded to tell everybody's fortune. The company gathered around her, and at the same time a girl came in and announced that Sylvia had sprained her ankle and couldn't come. Everybody said warmly, <Oh, how sickening, poor dear!> and forgot the subject immediately.

<We'll scoot off,> said Marjorie. <Never mind about saying good-bye. Nobody marks you. It's good luck about Sylvia, because she'll be at home and can't escape us. I sometimes wish they'd all sprain their ankles. And yet, you know, nearly all those people are doing very good work. Even the Kropotky crowd. I used to enjoy this kind of thing myself, once.>

<We're getting old, you and I,> said Wimsey. <Sorry, that's rude. But do you know, I'm getting on for forty, Marjorie.>

<You wear well. But you are looking a bit fagged tonight, Peter dear. What's the matter?>

<Nothing at all but middle-age.>

<You'll be settling down if you're not careful.>

<Oh, I've been settled for years.>

<With Bunter and the books. I envy you sometimes, Peter.>

Wimsey said nothing. Marjorie looked at him almost in alarm, and tucked her arm in his.

<Peter—do please be happy. I mean, you've always been the comfortable sort of person that nothing could touch. Don't alter, will you?>

That was the second time Wimsey had been asked not to alter himself; the first time, the request had exalted him; this time, it terrified him. As the taxi lurched along the rainy Embankment, he felt for the first time the dull and angry helplessness which is the first warning stroke of the triumph of mutability. Like the poisoned Athulf in the \textit{Fool's Tragedy}, he could have cried, <Oh, I am changing, changing, fearfully changing.> Whether his present enterprise failed or succeeded, things would never be the same again. It was not that his heart would be broken by a disastrous love—he had outlived the luxurious agonies of youthful blood, and in this very freedom from illusion he recognised the loss of something. From now on, every hour of light-heartedness would be, not a prerogative but an achievement—one more axe or case-bottle or fowling-piece, rescued, Crusoe-fashion, from a sinking ship.

For the first time, too, he doubted his own power to carry through what he had undertaken. His personal feelings had been involved before this in his investigations, but they had never before clouded his mind. He was fumbling—grasping uncertainly here and there at fugitive and mocking possibilities. He asked questions at random, doubtful of his object, and the shortness of the time, which would once have stimulated, now frightened and confused him.

<I'm sorry, Marjorie,> he said, rousing himself, <I'm afraid I'm being damned dull. Oxygen-starvation, probably. D'you mind if we have the window down a bit? That's better. Give me good food and a little air to breathe and I will caper, goat-like, to a dishonourable old age. People will point me out, as I creep, bald and yellow and supported by discreet corsetry, into the night-clubs of my great-grandchildren, and they'll say, <Look, darling! that's the wicked Lord~Peter, celebrated for never having spoken a reasonable word for the last ninety-six years. He was the only aristocrat who escaped the guillotine in the revolution of 1960. We keep him as a pet for the children.> And I shall wag my head and display my up-to-date dentures and say, <Ah, ha! They don't have the fun we used to have in my young days, the poor, well-regulated creatures!>>

<There won't be any night-clubs then for you to creep into, if they're as disciplined as all that.>

<Oh, yes—nature will have her revenge. They will slink away from the Government Communal Games to play solitaire in catacombs over a bowl of unsterilised skim-milk. Is this the place?>

<Yes; I hope there's someone to let us in at the bottom, if Sylvia's bust her leg. Yes—I hear footsteps. Oh, it's you, Eiluned; how's Sylvia?>

<Pretty all right, only swelled up—the ankle, that is. Coming up?>

<Is she visible?>

<Yes, perfectly respectable.>

<Good, because I'm bringing Lord~Peter Wimsey up, too.>

<Oh,> said the girl. <How do you do? You detect things, don't you? Have you come for the body or anything?>

<Lord~Peter's looking into Harriet Vane's business for her.>

<Is he? That's good. Glad somebody's doing something about it.> She was a short, stout girl with a pugnacious nose and a twinkle. <What do you say it was? I say he did it himself. He was the self-pitying sort, you know. Hullo, Syl—here's Marjorie, with a bloke who's going to get Harriet out of jug.>

<Produce him instantly!> was the reply from within. The door opened upon a small bed-sitting room, furnished with the severest simplicity, and inhabited by a pale, spectacled young woman in a Morris chair, her bandaged foot stretched out upon a packing-case.

<I can't get up, because, as Jenny Wren said, my back's bad and my leg's queer. Who's the champion, Marjorie?>

Wimsey was introduced, and Eiluned Price immediately inquired, rather truculently:

<Can he drink coffee, Marjorie? Or does he require masculine refreshment?>

<He's perfectly godly, righteous and sober, and drinks anything but cocoa and fizzy lemonade.>

<Oh! I only asked because some of your male belongings need stimulating, and we haven't got the wherewithal, and the pub's just closing.>

She stumped over to a cupboard, and Sylvia said:

<Don't mind Eiluned; she likes to treat 'em rough. Tell me, Lord~Peter, have you found any clues or anything?>

<I don't know,> said Wimsey. <I've put a few ferrets down a few holes. I hope something may come up the other end.>

<Have you seen the cousin yet—the Urquhart creature?>

<Got an appointment with him for tomorrow. Why?>

<Sylvia's theory is that he did it,> said Eiluned.

<That's interesting. Why?>

<Female intuition,> said Eiluned, bluntly. <She doesn't like the way he does his hair.>

<I only said he was too sleek to be true,> protested Sylvia. <And who else could it have been? I'm sure it wasn't Ryland Vaughan; he's an obnoxious ass, but he is genuinely heart-broken about it all.>

Eiluned sniffed scornfully, and departed to fill a kettle at a tap on the landing.

<And whatever Eiluned thinks, I can't believe Phil Boyes did it himself.>

<Why not?> asked Wimsey.

<He talked such a lot,> said Sylvia. <And he really had too high an opinion of himself. I don't think he would have wilfully deprived the world of the privilege of reading his books.>

<He would,> said Eiluned. <He'd do it out of spite, to make the grown-ups sorry. No, thanks,> as Wimsey advanced to carry the kettle, <I'm quite capable of carrying six pints of water.>

<Crushed again!> said Wimsey.

<Eiluned disapproves of conventional courtesies between the sexes,> said Marjorie.

<Very well,> replied Wimsey, amiably. <I will adopt an attitude of passive decoration. Have you any idea, Miss~Marriott, why this over-sleek solicitor should wish to make away with his cousin?>

<Not the faintest. I merely proceed on the old Sherlock Holmes basis, that when you have eliminated the impossible, then whatever remains, however improbable, must be true.>

<Dupin said that before Sherlock. I grant the conclusion, but in this case I question the premises. No sugar, thank you.>

<I thought all men liked to make their coffee into syrup.>

<Yes, but then I am very unusual. Haven't you noticed it?>

<I haven't had much time to observe you, but I'll count the coffee as a point in your favour.>

<Thanks frightfully. I say—can you people tell me just what was Miss~Vane's reaction to the murder?>

<Well\longdash> Sylvia considered a moment. <When he died—she was upset, of course\longdash>

<She was startled,> said Miss~Price, <but it's my opinion she was thankful to be rid of him. And no wonder. Selfish beast! He'd made use of her and nagged her to death for a year and insulted her at the end. And he was one of your greedy sort that wouldn't let go. She \textit{was} glad, Sylvia—what's the good of denying it?>

<Yes, perhaps. It was a relief to know he was finished with. But she didn't know then that he'd been murdered.>

<No. The murder spoilt it a bit—if it was a murder, which I don't believe. Philip Boyes was always determined to be a victim, and it was very irritating of him to succeed in the end. I believe that's what he did it for.>

<People do do that kind of thing,> said Wimsey, thoughtfully. <But it's difficult to prove. I mean, a jury is much more inclined to believe in some tangible sort of reason, like money. But I can't find any money in this case.>

Eiluned laughed.

<No, there never was much money, except what Harriet made. The ridiculous public didn't appreciate Phil Boyes. He couldn't forgive her that, you know.>

<Didn't it come in useful?>

<Of course, but he resented it all the same. She ought to have been ministering to his work, not making money for them both with her own independent trash. But that's men all over.>

<You haven't much opinion of us, what?>

<I've known too many borrowers,> said Eiluned Price, <and too many that wanted their hands held. All the same, the women are just as bad, or they wouldn't put up with it. Thank Heaven, I've never borrowed and never lent—except to women, and they pay back.>

<People who work hard usually do pay back, I fancy,> said Wimsey, <—except geniuses.>

<Women geniuses don't get coddled,> said Miss~Price, grimly, <so they learn not to expect it.>

<We're getting rather off the subject, aren't we?> said Marjorie.

<No,> replied Wimsey, <I'm getting a certain amount of light on the central figures in the problem—what journalists like to call the protagonists.> His mouth gave a wry little twist. <One gets a lot of illumination in that fierce light that beats upon a scaffold.>

<Don't say that,> pleaded Sylvia.

A telephone rang somewhere outside, and Eiluned Price went out to answer it.

<Eiluned's anti-man,> said Sylvia, <but she's a very reliable person.>

Wimsey nodded.

<But she's wrong about Phil—she couldn't stick him, naturally, and she's apt to think\longdash>

<It's for you, Lord~Peter,> said Eiluned, returning. <Fly at once—all is known. You're wanted by Scotland Yard.>

Wimsey hastened out.

<That you, Peter? I've been scouring London for you. We've found the pub.>

<Never!>

<Fact. And we're on the track of a packet of white powder.>

<Good God!>

<Can you run down first thing tomorrow? We may have it for you.>

<I will skip like a ram and hop like a high hill. We'll beat you yet, Mr~Bleeding Chief-Inspector Parker.>

<I hope you will,> said Parker, amiably, and rang off.

Wimsey pranced back into the room.

<Miss~Price's price has gone to odds on,> he announced. <It's suicide, fifty to one and no takers. I am going to grin like a dog and run about the city.>

<I'm sorry I can't join you,> said Sylvia Marriott, <but I'm glad if I'm wrong.>

<I'm glad I'm right,> said Eiluned Price, stolidly.

<And you are right and I am right and everything is quite all right,> said Wimsey.

Marjorie Phelps looked at him and said nothing. She suddenly felt as though something inside her had been put through a wringer.